 \begin{abstract}
Approximate caching offers a promising way to avoid repeated I/O and computation by serving vector search queries using results from past queries. Existing systems, however, have two main limitations. First, systems such as Proximity and Potluck use a single similarity threshold for all queries, ignoring local neighbor distributions and thus failing to simultaneously achieve high hit rate and recall. Second, all existing systems, including the state-of-the-art vector cache Aker, store each past query together with its top-$k$ results and reuse these results unchanged for similar queries. This frozen-answer design reduces recall and prevents cache hits from composing neighbors across multiple past queries, limiting both recall and hit-rate opportunities. Moreover, these \emph{query-materialized} caches must update cached top-$k$ lists upon insertions, since a newly inserted vector may belong to the top-$k$ results of multiple cached queries. In the worst case, this requires updating all cached entries, resulting in expensive updates or stale results.

We present QVCache, a \emph{data-materialized} approximate cache for vector search that addresses these limitations through two design decisions. First, an online learning algorithm maintains a distinct similarity threshold for each region of the embedding space and dynamically adapts it to the local neighbor distribution. Second, QVCache builds mini-indexes over cached data vectors and searches them to generate answers dynamically, without materializing query vectors or their top-$k$ results. This design enables cache hits to compose neighbors from multiple past queries and allows insertions to be handled efficiently by directly adding new vectors to the mini-indexes.

Integrating QVCache with pgvector, we evaluate it using the benchmark from prior work. QVCache achieves up to $X\%$ higher hit rate and $X\%$ higher recall than the state-of-the-art vector cache, Aker, while reducing average query latency by up to $Z\%$.

\end{abstract}

